\subsection{Три случая интегрирования дифференциальных биномов $x^n(a+bx^m)^p$}

\subsubsection*{1. Определение}

\textbf{Дифференциальным биномом} называется интеграл вида
\[
\int x^n\,(a + bx^m)^p\,dx,
\]
где $n, m, p \in \Q$, $a, b \in \R$, $b \neq 0$.

\begin{theorem}[Теорема Чебышёва]
Интеграл от дифференциального бинома выражается в элементарных функциях тогда и только тогда, когда хотя бы одна из трёх величин
\[ p, \quad \frac{n+1}{m}, \quad \frac{n+1}{m}+p \]
является целым числом.
\end{theorem}

\subsubsection*{2. Случай 1: $p \in \Z$}

Если $p$ — целое число, то замена $x = t^s$, где $s$ — общий знаменатель дробей $n$ и $m$, рационализирует интеграл.

\textbf{Пример.} $\displaystyle\int x^{1/2}(1+x^{1/3})^{-1}\,dx$: берём $s = 6$, $x = t^6$, $dx = 6t^5\,dt$:
\[
\int t^3\cdot\frac{1}{1+t^2}\cdot 6t^5\,dt = 6\int\frac{t^8}{1+t^2}\,dt.
\]

\subsubsection*{3. Случай 2: $\dfrac{n+1}{m} \in \Z$}

Обозначим $k = \dfrac{n+1}{m}$.

Делаем замену $t = a + bx^m$, то есть $x^m = \dfrac{t-a}{b}$:
\[
x = \left(\frac{t-a}{b}\right)^{1/m}, \qquad
dx = \frac{1}{m}\left(\frac{t-a}{b}\right)^{1/m - 1}\cdot\frac{dt}{b}.
\]

После замены $x^n\,dx$ содержит степень $(t-a)^{k-1}$ — целую степень, и интеграл становится рациональным по $t$.

\textbf{Пример.} $\displaystyle\int x^2(1+x^3)^{1/2}\,dx$: $n=2$, $m=3$, $\dfrac{n+1}{m} = 1 \in \Z$.

Замена $t = 1+x^3$, $dt = 3x^2\,dx$:
\[
\int t^{1/2}\cdot\frac{dt}{3} = \frac{1}{3}\cdot\frac{2}{3}t^{3/2}+C = \frac{2}{9}(1+x^3)^{3/2}+C.
\]

\subsubsection*{4. Случай 3: $\dfrac{n+1}{m} + p \in \Z$}

Обозначим $k = \dfrac{n+1}{m}+p$.

Делаем замену $t = \left(\dfrac{a}{x^m}+b\right)^{1/m}$, то есть $t^m = \dfrac{a}{x^m}+b$, что эквивалентно замене $x = \left(\dfrac{a}{t^m-b}\right)^{1/m}$.

После преобразований интеграл также становится рациональным по $t$.

\textbf{Пример.} $\displaystyle\int \frac{dx}{x^2\sqrt{1+x^2}}$: $n=-2$, $m=2$, $p=-\tfrac{1}{2}$.

Проверим: $\dfrac{n+1}{m}+p = \dfrac{-1}{2} - \dfrac{1}{2} = -1 \in \Z$.

Замена $t = \sqrt{\dfrac{1}{x^2}+1} = \dfrac{\sqrt{1+x^2}}{|x|}$, то есть $t^2 = \dfrac{1+x^2}{x^2}$:
\[
x = \frac{1}{\sqrt{t^2-1}}, \quad dx = -\frac{t}{(t^2-1)^{3/2}}\,dt, \quad
\sqrt{1+x^2} = \frac{t}{\sqrt{t^2-1}}.
\]
\[
\int \frac{-t\,dt/(t^2-1)^{3/2}}{1/(t^2-1)\cdot t/\sqrt{t^2-1}}
= -\int dt = -t + C = -\frac{\sqrt{1+x^2}}{|x|}+C.
\]

\subsubsection*{5. Сводная таблица}

\begin{center}
\renewcommand{\arraystretch}{1.8}
\begin{tabular}{|c|c|l|}
\hline
\textbf{Случай} & \textbf{Условие} & \textbf{Замена} \\
\hline
1 & $p \in \Z$ & $x = t^s$, $s = \mathrm{lcm}(\text{зн.}(n),\text{зн.}(m))$ \\
2 & $\dfrac{n+1}{m} \in \Z$ & $t = a + bx^m$ \\
3 & $\dfrac{n+1}{m}+p \in \Z$ & $t = \left(\dfrac{a}{x^m}+b\right)^{1/m}$ \\
\hline
\end{tabular}
\end{center}